\chapter{Finite detector access to a sharp physical minimum}
\label{ch:dg-finite-access}

An ideal observation law can determine a physical minimum while a finite
apparatus record remains unable to bound that minimum from above. The
distinction can be quantified in the same circle family. At fixed exposure
and positive read-noise variance, passive stationary trials and trials with
an oriented phase preparation have different sample requirements for a
specified binary comparison. Allowing read precision to consume a declared
measurement energy changes that comparison again. This chapter proves all
three statements in one explicit detector model.

The target is the already attained minimum compatible with the ideal
cosine-circle law,
\begin{equation}
 \mathcal M(\omega)=\frac{\omega^2}{\gamma},\qquad \gamma>0.
 \label{eq:dg-minimum-target}
\end{equation}
It is not an upper bound on the actual entropy of every hidden realization:
an invisible independently dissipating subsystem would already rule out
that stronger interpretation under perfect access. Here the detector
obstructs inference of a functional that is sharply fixed by the ideal
observed law itself.

\section{One fixed exposure and its moments}

Let the target diffusion satisfy
\[
 d\phi_t=\omega\,dt+\sqrt{2\gamma}\,dW_t\pmod{2\pi}.
\]
An exposure of fixed length $\tau>0$ yields
\begin{equation}
 A_\omega=\frac1\tau\int_0^\tau\cos\phi_t\,dt,
 \qquad X_\omega=A_\omega+\xi,
 \qquad \xi\sim N(0,\nu^2),\quad\nu>0,
 \label{eq:dg-exposure-record}
\end{equation}
where the read noise is independent of the entire target trajectory.
The parameters $\gamma$, $\tau$, and $\nu$ are known and fixed unless a
different precision law is explicitly declared below. In every case
$|A_\omega|\leq1$.

For passive stationary trials, the initial phase is uniform and
$\mathbb E A_\omega=0$. Define $V(\omega)=\mathbb E A_\omega^2$.
The covariance
$\mathbb E[\cos\phi_0\cos\phi_t]
 =\tfrac12e^{-\gamma|t|}\cos(\omega t)$ gives
\begin{equation}
 V(\omega)=\frac1{\tau^2}\operatorname{Re}
 \left\{\frac\tau z-\frac{1-e^{-z\tau}}{z^2}\right\},
 \qquad z=\gamma-i\omega.
 \label{eq:dg-exposure-variance}
\end{equation}
This follows by integrating the covariance over the exposure square, or
equivalently by integrating
$\tau^{-2}(\tau-t)e^{-\gamma t}\cos(\omega t)$ over $0\leq t\leq\tau$.

For an explicit asymptotic comparison, use the resonant sequence
\begin{equation}
 \omega_m=2\pi m/\tau,\quad m=1,2,\ldots,
 \qquad a_0=1-e^{-\gamma\tau}>0.
 \label{eq:dg-resonant-family}
\end{equation}
At these frequencies,
\begin{equation}
 V(\omega)=
 \frac{(\gamma\tau+a_0)\omega^2
       +(\gamma\tau-a_0)\gamma^2}
      {\tau^2(\gamma^2+\omega^2)^2}
 \sim\frac{\gamma\tau+a_0}{\tau^2\omega^2}.
 \label{eq:dg-resonant-variance}
\end{equation}
In particular,
\begin{equation}
 \Delta_V(\omega):=|V(\omega)-V(2\omega)|
 \sim\frac{3(\gamma\tau+a_0)}{4\tau^2\omega^2}.
 \label{eq:dg-variance-gap}
\end{equation}
All asymptotic statements in this chapter fix the detector and confidence
constants and take $m\to\infty$ along this sequence.

Two moment bounds hold uniformly over the initial phase distribution,
and hence apply to both stationary and prepared trials. Set
\begin{align}
 H_2&=\frac{[2+\gamma\tau+\sqrt{2\gamma\tau}]^2}{\tau^2},\\
 H_4&=\frac{8[(2+\gamma\tau)^4+12\gamma^2\tau^2]}{\tau^4}.
 \label{eq:dg-ito-constants}
\end{align}
For $\omega\ne0$,
\begin{equation}
 \mathbb E A_\omega^2\leq\frac{H_2}{\omega^2},
 \qquad
 \mathbb E A_\omega^4\leq\frac{H_4}{\omega^4}.
 \label{eq:dg-ito-moments}
\end{equation}
To see this, It\^o's formula gives
\begin{equation}
 \omega\tau A_\omega
 =\sin\phi_\tau-\sin\phi_0
   +\gamma\int_0^\tau\sin\phi_s\,ds
   -\sqrt{2\gamma}\int_0^\tau\cos\phi_s\,dW_s.
 \label{eq:dg-ito-exposure}
\end{equation}
The bounded part has absolute value at most $2+\gamma\tau$. For
$M_t=\int_0^t\cos\phi_s\,dW_s$, It\^o isometry gives
$\mathbb E M_\tau^2\leq\tau$ and the fourth-moment identity gives
\[
 \mathbb E M_\tau^4
 =6\mathbb E\int_0^\tau M_s^2\cos^2\phi_s\,ds
 \leq6\int_0^\tau s\,ds=3\tau^2.
\]
Minkowski's inequality proves the second-moment bound. The elementary
inequality $(x+y)^4\leq8(x^4+y^4)$ proves the fourth-moment bound.
No independence between the bounded and martingale terms is assumed.

\section{A finite record cannot provide a uniformly finite upper bound}

Fix $n$ predetermined exposure windows, each of width $\tau$, from a
stationary trajectory. The windows may overlap and their latent exposure
values may be correlated. Let $P_\omega^{(n)}$ be the law of their noisy
records under independent Gaussian read errors, and let
$Q^{(n)}=N(0,\nu^2I_n)$. Convexity of relative entropy under Gaussian
mixing yields
\begin{equation}
 D(P_\omega^{(n)}\Vert Q^{(n)})
 \leq\frac{\mathbb E\|A_\omega^{(n)}\|^2}{2\nu^2}
 =\frac{nV(\omega)}{2\nu^2}
 \leq\frac{nH_2}{2\nu^2\omega^2}.
 \label{eq:dg-mixture-kl}
\end{equation}
Indeed, conditional on the exposure vector, the readout is Gaussian with
that vector as its mean and covariance $\nu^2I_n$. Its relative entropy
to $Q^{(n)}$ is the squared mean divided by $2\nu^2$; mixing can only
decrease the relative entropy. Pinsker's inequality therefore implies
$\operatorname{TV}(P_\omega^{(n)},Q^{(n)})\to0$ as
$\omega\to\infty$, whereas $\mathcal M(\omega)\to\infty$.

\begin{theorem}[Upper-confidence obstruction for the sharp minimum]
\label{thm:dg-upper-obstruction}
Suppose a measurable extended-valued statistic $U$ of this finite record
satisfies
\[
 P_\omega^{(n)}\{U\geq\mathcal M(\omega)\}\geq1-\alpha
 \quad\hbox{for every }\omega>0,
 \qquad 0<\alpha<1.
\]
Then $Q^{(n)}\{U=+\infty\}\geq1-\alpha$. In particular, no such
uniformly covered upper bound is finite almost surely under even one of the
finite-$\omega$ record laws. The conclusion concerns the fixed finite
detector design and unbounded frequency class just stated.
\end{theorem}

\begin{proof}
For every finite $K$, sufficiently large $\omega$ has
$\mathcal M(\omega)>K$, so
\[
 Q^{(n)}\{U>K\}
 \geq1-\alpha-\operatorname{TV}(P_\omega^{(n)},Q^{(n)}).
\]
Let $\omega\to\infty$ and then $K\to\infty$ to obtain the first
assertion. Each Gaussian mixture $P_\omega^{(n)}$ has a strictly positive
density everywhere on $\mathbb R^n$, and is therefore mutually absolutely
continuous with $Q^{(n)}$. If $U$ were finite almost surely under any one
$P_\omega^{(n)}$, it would also be finite almost surely under $Q^{(n)}$,
contradicting the first assertion.
\end{proof}

The pure-noise law is a limiting reference, not necessarily a finite-parameter
member of the model class. Accordingly, the theorem does not claim
$P_\omega^{(n)}\{U=+\infty\}\geq1-\alpha$ at every fixed $\omega$.
The general quantitative statement is
\[
 P_\omega^{(n)}\{U=+\infty\}
 \geq1-\alpha-\operatorname{TV}(P_\omega^{(n)},Q^{(n)}).
\]
A calibrated frequency or generator envelope removes the unbounded class
used in this argument. Adaptive timing, adaptive precision, or a growing
number of measurements needs its own sequential analysis; it is not included
in this fixed-design theorem.

\section{Matching passive sample requirements at fixed precision}

Now compare the known alternatives $\omega$ and $2\omega$ using
independent stationary trials, each with a fresh stationary preparation and
one exposure. Define $n^*_{\rm passive}(\omega,\epsilon)$ as the minimum
number of such trials needed for equal-prior binary error at most
$\epsilon\in(0,1/2)$. An upper construction below bounds each of the two
conditional errors by $\epsilon$, which is a stronger guarantee.

The zero stationary mean improves the mixture bound for independent trials.
Let $P$ be the one-record law of $A+\xi$, $Q=N(0,\nu^2)$, and let
$A'$ be an independent copy of $A$. Direct multiplication of Gaussian
likelihood ratios gives
\begin{equation}
 \chi^2(P\Vert Q)
 =\mathbb E e^{AA'/\nu^2}-1
 =\sum_{k=1}^\infty
       \frac{(\mathbb E A^k)^2}{k!\nu^{2k}}.
 \label{eq:dg-mixture-chi-square}
\end{equation}
Bounded $A$ justifies termwise expectation. Since $\mathbb E A=0$ and
$|\mathbb E A^k|\leq\mathbb E A^2=V$ for every $k\geq2$,
\begin{equation}
\chi^2(P\Vert Q)\leq C_{\nu,0} V^2,
 \qquad C_{\nu,0}=e^{1/\nu^2}-1-1/\nu^2>0.
 \label{eq:dg-stationary-chi-square}
\end{equation}
For the stationary circle the bound is stronger. Shifting every phase by
$\pi$ preserves its stationary path law and changes $A$ to $-A$.
All odd moments therefore vanish, and retaining only the even terms in
Equation~\eqref{eq:dg-mixture-chi-square} gives
\begin{equation}
 \chi^2(P\Vert Q)\leq C_\nu V^2,
 \qquad C_\nu=\cosh(1/\nu^2)-1.
 \label{eq:dg-symmetric-chi-square}
\end{equation}
The sharper circle constant is used in the count bounds below.
For independent trials, relative entropy tensorizes. Using
$D\leq\log(1+\chi^2)\leq\chi^2$ and the triangle inequality for total
variation,
\[
 \operatorname{TV}(P_\omega^{\otimes n},P_{2\omega}^{\otimes n})
 \leq\sqrt{\frac{nC_\nu}{2}}\{V(\omega)+V(2\omega)\}.
\]
The optimal equal-prior error is $(1-\operatorname{TV})/2$. Hence a necessary
sample count is
\begin{equation}
 n\geq
 \frac{2(1-2\epsilon)^2}
      {C_\nu\{V(\omega)+V(2\omega)\}^2}.
 \label{eq:dg-passive-necessary}
\end{equation}
The constant can be numerically weak, especially for small $\nu$; a bound
below one is vacuous as an integer-count constraint. The asymptotic order
remains informative at fixed positive $\nu$.

For a matching upper bound, use the sample mean of $X^2$ and put its threshold
midway between $\nu^2+V(\omega)$ and $\nu^2+V(2\omega)$. Gaussian moments
and $|A|\leq1$ give
\begin{align}
 \operatorname{Var}(X^2)
 &=\mathbb E A^4-V^2+4\nu^2V+2\nu^4\notag\\
 &\leq S_\nu(V):=(1+4\nu^2)V-V^2+2\nu^4.
 \label{eq:dg-square-variance}
\end{align}
Chebyshev's inequality then bounds each conditional error by $\epsilon$
whenever
\begin{equation}
 n\geq
 \left\lceil
 \frac{4\max\{S_\nu(V(\omega)),S_\nu(V(2\omega))\}}
      {\epsilon\,\Delta_V(\omega)^2}
 \right\rceil.
 \label{eq:dg-passive-sufficient}
\end{equation}
A simpler valid numerator is $4(1+6\nu^2+3\nu^4)$, obtained by bounding
$\mathbb E X^4$ directly. The more precise expression above improves
constants without changing the access or confidence criterion.

Equations~\eqref{eq:dg-resonant-variance}--\eqref{eq:dg-variance-gap}
and the two sample bounds prove
\begin{equation}
 n^*_{\rm passive}(\omega_m,\epsilon)=\Theta(\omega_m^4)
      =\Theta(\mathcal M(\omega_m)^2).
 \label{eq:dg-passive-rate}
\end{equation}
The constants may depend on fixed $\gamma$, $\tau$, $\nu$, and
$\epsilon$. This is a binary comparison along the declared sequence,
not a minimax theorem for estimating an arbitrary unknown frequency from
an arbitrary trajectory. Nonoverlapping windows from one stationary
trajectory are not automatically independent stationary preparations.

\section{Matching requirements with an oriented phase preparation}

Suppose each independent trial can instead start from the actual known
phase $\phi_0=\pi/2$, under both hypotheses. Its expectation during the
exposure is $\mathbb E\cos\phi_t=-e^{-\gamma t}\sin(\omega t)$.
At resonance, the one-record mean is
\begin{equation}
 \mu(\omega)=-\frac{a_0\omega}{\tau(\gamma^2+\omega^2)},\qquad
 \Delta_\mu(\omega)
 =\frac{a_0|\omega|\,|2\omega^2-\gamma^2|}
        {\tau(\gamma^2+\omega^2)(\gamma^2+4\omega^2)}
 \sim\frac{a_0}{2\tau|\omega|}.
 \label{eq:dg-prepared-mean}
\end{equation}
The gap is positive for sufficiently large resonant frequencies; at
$\omega^2=\gamma^2/2$ this particular mean test has zero gap and supplies
no discrimination guarantee.

Since $A\in[-1,1]$, its centered version is sub-Gaussian with variance proxy
one. Adding the independent Gaussian read noise gives proxy $1+\nu^2$.
The sample-mean midpoint test therefore bounds each conditional error by
$\epsilon$ with the sufficient count
\begin{equation}
 n\geq
 \left\lceil
 \frac{8(1+\nu^2)\log(1/\epsilon)}{\Delta_\mu(\omega)^2}
 \right\rceil.
 \label{eq:dg-prepared-sufficient}
\end{equation}

A lower bound of the same order follows without assuming stationarity.
The uniform second-moment bound~\eqref{eq:dg-ito-moments} and the
Gaussian-mixture relative-entropy inequality give, for any common fixed
phase-preparation rule,
\[
 D(P_\omega^{\otimes n}\Vert Q^{\otimes n})
 \leq\frac{nH_2}{2\nu^2\omega^2}.
\]
For the pair $\omega,2\omega$, Pinsker and the triangle inequality imply
\[
 \operatorname{TV}(P_\omega^{\otimes n},P_{2\omega}^{\otimes n})
 \leq\frac{3\sqrt{nH_2}}{4\nu|\omega|}.
\]
An equal-prior error at most $\epsilon$ consequently requires
\begin{equation}
 n\geq\frac{16\nu^2(1-2\epsilon)^2\omega^2}{9H_2}.
 \label{eq:dg-prepared-necessary}
\end{equation}
The construction and obstruction together prove
\begin{equation}
 n^*_{\rm prepared}(\omega_m,\epsilon)=\Theta(\omega_m^2)
         =\Theta(\mathcal M(\omega_m)).
 \label{eq:dg-prepared-rate}
\end{equation}

The preparation need not be a perfect point mass. A known initial phase
law with $\mathbb E\cos\phi_0=0$ and a fixed nonzero
$r=\mathbb E\sin\phi_0$ multiplies the mean in
Equation~\eqref{eq:dg-prepared-mean} by $r$, preserving the exponent with
a sufficient-count constant proportional to $r^{-2}$. Merely preparing
$Y_0=\cos\phi_0=0$ does not establish this condition: an equal mixture of
$\pi/2$ and $3\pi/2$ has $r=0$ and cancels the mean signal.

These prepared records are not stationary entropy-production records and are
not inserted into the stationary local-defect certificate. They distinguish
two generators whose stationary ideal laws have the target minima
$\mathcal M(\omega)$ and $4\mathcal M(\omega)$. A phase preparation is a
change in permitted access, with its own apparatus and energy requirements.

\section{What changes when read precision consumes energy}

The previous exponents count trials at fixed positive read-noise variance.
They are not total-energy exponents. To make that distinction explicit,
consider the following additional, declared sensor model:
\begin{equation}
 \nu_i^2=\frac{\kappa}{E_i},\qquad \kappa>0,
 \quad E_i>0,
 \label{eq:dg-sensor-energy-law}
\end{equation}
where $E_i$ is the measurement energy allocated to record $i$. This is an
assumed calibration law, not a universal thermodynamic law. The finite
energy schedule is fixed before the data, and read errors remain independent
of the target and each other. Preparation, reset, exposure duration, and
actuation costs are not included in $E_i$.

Let $E_{\rm meas}=\sum_iE_i$, and use the Gaussian reference with diagonal
variances $\kappa/E_i$. Conditional Gaussian relative entropy and
Equation~\eqref{eq:dg-ito-moments} imply
\begin{equation}
 D(P_\omega\Vert Q)
 \leq\sum_i\frac{E_i\mathbb E A_{\omega,i}^2}{2\kappa}
 \leq\frac{H_2E_{\rm meas}}{2\kappa\omega^2}.
 \label{eq:dg-energy-kl}
\end{equation}
This bound permits correlated latent exposures at predetermined windows,
provided each has the stated dynamics and duration. The initial phases can
follow the same fixed preparation rule under both hypotheses. The
corresponding pairwise Pinsker bound gives the necessary measurement energy
\begin{equation}
 \boxed{\displaystyle
 E_{\rm meas}\geq
 \frac{16\kappa(1-2\epsilon)^2\omega^2}{9H_2}.}
 \label{eq:dg-energy-necessary}
\end{equation}
The proof does not cover an energy schedule adaptively selected from noisy
records or an intervention that changes the dynamics during exposure.

Passive stationary trials can attain the same exponent by buying more
precision per record. Fix $c>0$ and use
$E_i=c\omega^2$ under both alternatives, so
$\nu_i^2=\kappa/(c\omega^2)$. This design is permitted because the
binary candidate pair is known. Let
\begin{equation}
 B_c=H_4+\frac{6\kappa H_2}{c}
             +\frac{3\kappa^2}{c^2}.
 \label{eq:dg-energy-fourth-constant}
\end{equation}
Under either hypothesis, the Gaussian fourth-moment expansion and
Equation~\eqref{eq:dg-ito-moments} give
\[
 \mathbb E X^4\leq\frac{B_c}{\omega^4}.
\]
The midpoint test of the sample mean of $X^2$ consequently has both errors
at most $\epsilon$ if
\begin{equation}
 n\geq
 \left\lceil\frac{4B_c}
       {\epsilon\,\omega^4\Delta_V(\omega)^2}\right\rceil.
 \label{eq:dg-energy-passive-sufficient}
\end{equation}
By Equation~\eqref{eq:dg-variance-gap}, the right-hand side remains
bounded along the high-frequency resonant sequence. Thus a bounded number
of passive trials, each with energy proportional to $\omega^2$, suffices.
Together with~\eqref{eq:dg-energy-necessary}, this establishes
\begin{equation}
 E^*_{\rm meas,passive}(\omega_m,\epsilon)=\Theta(\omega_m^2).
 \label{eq:dg-passive-energy-rate}
\end{equation}
Prepared trials at fixed $\nu$ already attain this measurement-energy
order, because each such readout has fixed energy $\kappa/\nu^2$ and
their required count is $\Theta(\omega^2)$. Therefore, under the declared
sensor law and allowing precision allocation,
\begin{equation}
 \begin{aligned}
 E^*_{\rm meas,prepared}(\omega_m,\epsilon)
   &=\Theta(\omega_m^2),\\
 E^*_{\rm meas,passive}(\omega_m,\epsilon)
   &=\Theta(\omega_m^2).
 \end{aligned}
 \label{eq:dg-energy-comparison}
\end{equation}
Both are of order $\mathcal M(\omega_m)$ at fixed $\gamma$; their optimal
constants and procedures have not been shown equal. Preparation does improve the fixed-precision
sample exponent; this model supplies no improvement in the optimized
measurement-energy exponent. A total-energy comparison would additionally
need calibrated preparation, reset, control, exposure, and timing costs.

\section{Operational scope and reproducible design}

At fixed precision, $n$ independent trials consume $n$ preparations,
$n$ exposed scalar readouts, and $n\tau$ units of declared evolution time.
The passive and prepared trial classes have different preparation access.
Under the precision-energy model, passive inference can instead use a
bounded trial count with increasing per-read energy. None of these
accountings identifies declared evolution time with total laboratory time
or with total energy.

The comparison closes a restricted operational problem: discrimination of
known candidate generators along a specified high-frequency family, with
fixed exposure and a declared observation channel. It is supported by exact
moment identities, Gaussian-mixture information bounds, and explicit tests.
It does not establish a universal data cap, an optimal estimator for every
unknown diffusion, or a new physical energy law. The elementary statistical
ingredients are established; any narrower novelty lies in the particular
physical target, access comparison, and assumptions, subject to prior-art
review. Numerical design tables evaluate these proved formulas and their
resource counts; they are not laboratory measurements or simulated detection
successes unless separately accompanied by such records.
